\section{A certifiable analytic learning lane}
\subsection{Targets, fields and arithmetic scope}
The output is a complex holomorphic function, not an unconstrained real MLP
on $\mathbb R^2$. Fix a center $c$ and a finite family of targets $F_\chi$
analytic on a neighborhood of the closed radius-$r$ disk. Certified values
of $M_{r,\chi}\ge\sup|F_\chi|$ and of an evaluation error $\eta$ are inputs
to the quantitative theorem. They must be obtained without GRH or the desired
arithmetic positivity. A finite conductor family uses the maximum of its
independently bounded constants; this is not a uniform theorem over all
conductors. Arithmetic character data are exact discrete inputs, not rounded
embedding labels.

For primitive nonprincipal Dirichlet characters, the completed function is
entire and a direct bound is given below. Pole-cleared zeta/Dedekind targets
and known holomorphic Hecke or elliptic targets require their own explicit
normalization and growth bounds. General Artin holomorphy is not assumed.
Approximating on a compact domain avoiding possible poles does not prove that
those poles are absent. Real/complex analytic norms, finite-field algebra and
tropical semirings must likewise retain their distinct meanings.

\subsection{A finite holomorphic readout theorem}
Write $F(c+z)=\sum_{n\ge0}a_nz^n$ and choose
$0<R<\rho<r$, integers $m>N\ge0$ and sample angles $\theta_j=2\pi j/m$.
Let $|y_j-F(c+\rho e^{i\theta_j})|\le\eta$, including all quadrature and
rounding errors. The model and empirical objective are
\[
 p_b(c+z)=\sum_{n=0}^N b_n(z/\rho)^n,\qquad
 \mathcal L(b)=\frac1{2m}\sum_{j=0}^{m-1}
 \left|\sum_{n=0}^Nb_ne^{in\theta_j}-y_j\right|^2.
\]
This is a fixed holomorphic feature layer and a learned complex linear
readout. It is a genuine convex learning problem; no nonlinear feature-learning
claim is attached to its guarantee.

\begin{proposition}[Four-part uniform error bound]\label{hl:finite}
Put $u=R/\rho$, $v=R/r$ and $w=(\rho/r)^m$. Let
$b_n^*=m^{-1}\sum_jy_je^{-in\theta_j}$, and take gradient descent with
$0<\gamma<2$. After $k$ steps from $b^0$, on $|s-c|\le R$,
\[
 \begin{split}
 |p_{b^k}(s)-F(s)|\le{}&
 \frac{|1-\gamma|^k\|b^0-b^*\|_2}{\sqrt{1-u^2}}
 +\eta\frac{1-u^{N+1}}{1-u}\\
 &+M_r\frac{w}{1-w}\frac{1-v^{N+1}}{1-v}
 +M_r\frac{v^{N+1}}{1-v}.
 \end{split}
\]
The terms are optimization, certified sample error, Cauchy aliasing and
analytic truncation, respectively.
\end{proposition}
\begin{proof}
Discrete Fourier orthogonality gives $A^*A/m=I$ for
$A_{jn}=e^{in\theta_j}$, since $m>N$. In real and imaginary coordinates the
gradient is $b-b^*$; hence
$b^k-b^*=(1-\gamma)^k(b^0-b^*)$. Cauchy--Schwarz bounds its evaluation by
$\|b^k-b^*\|_2(\sum_{n=0}^Nu^{2n})^{1/2}$, giving the first term.
The displayed iteration is in exact arithmetic. A numerical FFT/GD implementation
must also bound coefficient-rounding error; a coefficient error of norm $d_k$
adds at most $d_k/\sqrt{1-u^2}$ to the evaluation budget.
This is convergence to the projected least-squares optimum; the empirical
residual need not vanish when the labels have components outside the model.

For exact samples, the Fourier coefficient divided by $\rho^n$ is
$a_n+\sum_{\ell\ge1}a_{n+\ell m}\rho^{\ell m}$.
Uniform convergence on the sampling circle permits interchanging sums.
Cauchy's estimate $|a_j|\le M_rr^{-j}$ bounds the alias by
$M_rr^{-n}w/(1-w)$. Sample errors contribute at most $\eta\rho^{-n}$.
Multiply by $R^n$, sum the two finite geometric series and add
$\sum_{n>N}M_r(R/r)^n$. These are exactly the last three terms.
\end{proof}

Choose $N$, then $m>N$, then $\eta$ and the optimization tolerance so that
each displayed term is at most $\varepsilon/4$. All four can tend to zero;
thus any fixed compact subdisk and finite independently bounded target family
have arbitrary-accuracy approximation. This claim needs a certified oracle
bound, not an empirical test-set MSE. The present high-precision demonstrator
is explicitly numerical and does not relabel its floating evaluations as
interval certificates.

\subsection{The infinite-width limit and its actual rate}
Choose $R<\rho<r_0<r$ and use features $\phi_n(z)=(z/r_0)^n$.
Their kernel has the compact-uniform limit
\[
 K_N(z,w)=\sum_{n=0}^N\phi_n(z)\overline{\phi_n(w)}
 \longrightarrow K(z,w)=\frac1{1-z\bar w/r_0^2}.
\]
The geometric tail proves this limit directly; positivity follows because
each quadratic form is a sum of squared feature combinations. The readout
space is the sequence space $\ell^2$ with bounded evaluation on $|z|<r_0$.
On the radius-$\rho$ circle its population coefficient Gram operator is
diagonal with eigenvalues $\lambda_n=(\rho/r_0)^{2n}$.

\begin{proposition}[Infinite-feature readout convergence]\label{hl:infinite}
Suppose the initial coefficient error $e(0)$ is in $\ell^2$, as for zero
initialization and the analytic target above. Population gradient flow has
$e_n(t)=e^{-\lambda_nt}e_n(0)$ and converges in $\ell^2$, hence uniformly on
every $|z|\le R<r_0$. If $\beta>0$ and
$\sum_n\lambda_n^{-2\beta}|e_n(0)|^2\le B_\beta^2$, then
\[
 \|p_t-F\|_{|z|\le R}
 \le\frac{B_\beta}{\sqrt{1-(R/r_0)^2}}
       \left(\frac{\beta}{et}\right)^\beta,\qquad t>0.
\]
For zero initialization this source condition holds whenever
$(r_0/r)(r_0/\rho)^{2\beta}<1$, with
$B_\beta\le M_r[1-((r_0/r)(r_0/\rho)^{2\beta})^2]^{-1/2}$.
\end{proposition}
\begin{proof}
The diagonal loss gives the scalar ODE in each coordinate. Dominated
convergence, with majorant $|e_n(0)|^2$, proves norm convergence. Evaluation
is bounded by the geometric-series factor. Under the source condition,
$\|e(t)\|_2\le B_\beta\sup_{\lambda>0}\lambda^\beta e^{-t\lambda}$;
the maximum occurs at $\lambda=\beta/t$. Cauchy's coefficient estimate gives
$|a_n r_0^n|\le M_r(r_0/r)^n$, and summing the resulting geometric series
proves the explicit source bound.
\end{proof}

There is no width-uniform exponential rate: $\lambda_n\to0$. At finite
degree the slowest mode has rate $(\rho/r_0)^{2N}$. Fourier whitening recovers
the identity Gram matrix in Proposition~\ref{hl:finite}, but its unbounded
infinite-dimensional preconditioner must not be called a free uniform rate.
Ridge adds a positive rate floor and a calculable bias; its coefficient bias
factor is $\alpha/(\lambda_n+\alpha)$, and $\alpha$ must tend to zero for
unbiased convergence. A standard Gaussian-process initialization need not
lie in the kernel's RKHS; the stated $\ell^2$ hypothesis cannot be silently
applied to it. These exact fixed-feature statements are distinct from the
nonlinear NTK limits of deep networks \cite{HLNTK} and from the trained
few-qubit feature-learning circuits.

\subsection{An unconditional primitive-Dirichlet size oracle}
For a primitive nonprincipal character of conductor $q$, parity $a\in\{0,1\}$
and root number $|\epsilon_\chi|=1$, put $d=\pi/q$ and
$\theta_\chi(t)=\sum_{n\ge1}\chi(n)n^a e^{-dn^2t}$.
The standard primitive theta transformation and Mellin identity give
\[
 \Lambda(s,\chi)=\int_1^\infty\theta_\chi(t)t^{(s+a)/2-1}\,dt
 +\epsilon_\chi\int_1^\infty\theta_{\bar\chi}(t)t^{(1-s+a)/2-1}\,dt.
\]
The character conjugation and root number are retained. These standard
functional/analytic inputs are unconditional in the stated Dirichlet scope
\cite{HLDirichlet}; this text does not claim a new Lean proof of Poisson
summation. Exponential decay makes both integrals entire in $s$.

For $|s-c|\le r$, choose an integer
$k\ge\max(0,(1+|\operatorname{Re}c|+r+a)/2-1)$ and define
\[
 Q_k(x)=\sum_{j=0}^k\frac{k!}{(k-j)!x^{j+1}}.
\]
Integration by parts gives
$\int_1^\infty t^ke^{-xt}\,dt=e^{-x}Q_k(x)$.
Consequently $M_r=2\sum_{n\ge1}n^ae^{-dn^2}Q_k(dn^2)$ is a valid
finite upper bound. For a computable finite version, choose $n_0\ge1$ and sum through $n_0-1$;
the remaining one-integral sum is at most
\[
 C_a Q_k(dn_0^2)
 \frac{e^{-dn_0^2/2}}{1-e^{-dn_0}},\qquad
 C_0=1,\quad C_1=(de)^{-1/2}.
\]
Indeed $n^ae^{-dn^2/2}\le C_a$, $Q_k$ decreases, and
$(n_0+j)^2\ge n_0^2+2n_0j$. Twice this tail covers both Mellin terms.
The same bound controls theta truncation, with no assumption about zeros.

After truncating the $n$ sum and the $t$ integral, elementary interval
exponential, logarithm, sine and cosine evaluation on a compact interval,
with a bounded-derivative quadrature remainder, gives an arbitrary requested
certified $\eta$. The explicit exponential tails control both truncations.
This proves computability of the oracle; it does not assert that an ordinary
high-precision special-function call is already a validated implementation.
For a finite conductor cutoff, take the maximum of the computed bounds.

\subsection{Additional objectives and what they can certify}
\subsubsection{Contour-aware approximation and active learning.}
Augment the coefficient objective with independent interval enclosures of
$|p-F|$ on contours, derivative-weighted losses and adaptive samples near
small $|p|$. If $\sup_C|p-F|<\inf_C|p|$ on a Jordan contour, Rouch\'e's
theorem transfers the exact interior zero count of $p$ to $F$. An interval
root count and a strict margin are acceptance conditions. Sampling finitely
many heights alone proves no exhaustion of the off-critical domain.

\subsubsection{Arithmetic and functional-equation residuals.}
Use exact character masks, conductor dependence, primitive Euler factors,
gamma normalization and the conjugate functional equation as structural
features or independently verified residuals. On $\operatorname{Re}s\ge1+\delta$,
with $\delta>0$, for $N\ge3$ the logarithmic-derivative Dirichlet tail has the elementary bound
$N^{-\delta}(\log N/\delta+1/\delta^2)$, from
$\Lambda(n)\le\log n$ and integral comparison. This is a useful certified
training region; it does not extend absolute convergence into the critical
strip. A functional equation alone also does not locate zeros.

\subsubsection{Operator and determinant refinement.}
Learn compatible finite compressions, graph-norm/domain residuals,
self-adjoint boundary conditions and uniform Schatten-tail bounds from the
actual infinite dessin/order construction. Trace-norm convergence transfers
ordinary Fredholm determinants, with the standard bound
\[
 |\det(I+A)-\det(I+B)|
 \le\|A-B\|_1e^{1+\|A\|_1+\|B\|_1}
\]
\cite{HLFredholm}. Strong resolvent convergence alone does not supply this
bound. A reciprocal critical-spectrum model typically needs a genus-one
regularized determinant and Hilbert--Schmidt control; it must not silently
be treated as trace class. The arithmetic determinant identification remains
an independent obligation even when the learned operators are Hermitian.

\subsubsection{Certificate-guided moment and algebra learning.}
For the actual original-family low-gcd moment, use adversarial legal masks,
conductors, short rows and retained II--I covariance to search for a sharp
envelope or counterexample. Rational/interval or exact symbolic checking must
certify a proposed inequality uniformly. A sum-of-squares or dual certificate
can prove its stated parameter box; extrapolating it to all conductors needs
a separate analytic argument. Preserve exact locally gentle relations,
surface-order completions, determinant semi-invariant weights and module
category labels in the representation architecture. Small relation loss is
not an exact functor or an arithmetic comparison.

\subsubsection{Quantum kernels, width and proof search.}
Quantum feature Gram positivity is not Weil/arithmetic positivity. A proved
finite Gram lower bound gives a readout rate, while growing width requires
the actual limiting kernel and initialization hypotheses. Proof-guided
co-training can learn proposed identities or compression maps, then export
rational/algebraic candidates for independent Lean/interval validation.
Validation, rather than a confidence score, decides which claims survive.

\subsection{A decisive obstruction to a false universal guarantee}
\begin{proposition}\label{hl:real-root-floor}
No sequence of real-rooted polynomials converges locally uniformly on
$\mathbb C$ to $1+z^2$. Every such polynomial has error at least $7/16$
somewhere on $|z-i|=1/4$.
\end{proposition}
\begin{proof}
On that circle, $|1+z^2|=|z-i||z+i|\ge(1/4)(2-1/4)=7/16$.
A smaller uniform error would, by Rouch\'e, force a zero of the polynomial
inside the disk, whereas all of its zeros are real and the disk misses the
real axis. The same argument excludes local-uniform convergence.
\end{proof}
Thus a claim that an everywhere-real-zero/self-adjoint determinant family
universally approximates the completed arithmetic target may already contain
the missing GRH bridge. Hurwitz transfers zero-freeness only after independent
local-uniform convergence and a nonzero limit have been proved. The analytic
learner above deliberately proves approximation without asserting that
spectral restriction. Its learned functions are evidence-bearing candidates,
not a completed GRH proof. All new analytic/probabilistic statements here are
prose proofs and scoped computations; no complete new Lean layer is claimed.
